\scanpage{81}
\mspage{30}

Comme les hyp. sont stables par chang. de base, il suffit de prouver $f\in W$. Supposons p. ex. $f$ propre, prouvons que $f$ est $W$-\uline{asphérique}, i.e. que les $X_{/y}$ ($y\in\mathrm{Ob}\,Y$) sont $W$-asphériques. Mais $X_y\longrightarrow X_{/y}$ étant $W$-coasphérique par hyp. de propreté, est $\in W$ (par $W(2\,\mathrm{bis})$) et comme $X_y$ est $W$-asph. par hyp., $X_{/y}$ aussi. Mais par $W(2\,\mathrm{bis})$, $f$ $W$-asphérique $\Longrightarrow f\in W$. Le cas $f$ lisse est dual.

\deleted{\uline{Remarque}} Ainsi, \deleted{\unsure{si}} $f$ est propre, s'il est à fibres $W$-asphériques il est ``universellement $W$-asphérique'', i.e. est et reste $W$-asphérique par \deleted{\unsure{tout}} ch. \deleted{\unsure{gement}} de base. Inversement, si $f$ est univ. \inserted{dans} $W$\deleted{-asphérique} (pas forcément propre) alors \inserted{\unsure{en particulier|en vertu de}} ses fibres sont $W$-asphériques \deleted{Cela résulte du fait que \deleted{\unsure{$\xrightarrow{f}$}} $Z\longrightarrow e$ est $W$-asphérique (ou $W$-coasphérique) \unsure{sss|ssi} $Z$ est $W$-asphérique (trivial). Notons que $f$ peut être universellement $W$-asphérique c'est \unsure{-- à la fois univ. $W$-asphér.}} Mais 5. n'implique bien sûr pas que $f$ soit propre, ou lisse.

% slanted
\margin{Mais ``universellt dans $W$'' implique déjà ``universellement $W$-asphérique'', \unsure{NB} la notion de $W$-asphéricité est définie par prop. 1 a).}

\uline{Corollaire} Si $W$ satisfait $W(2)$ (pas seulement $W(2\,\mathrm{bis})$), alors pour un $S$-morphisme $X\longrightarrow Y$ de cat. qui sont toutes deux \deleted{\unsure{$W$-}}propres sur $S$, ou \deleted{\unsure{toutes}} deux $W$-lisses, si $f_s:X_s\mathrel{\relbar}Y_s$ est dans $W$ pour tt $s\in S$, alors $f\in W$.

% archive sheet number: 81
